\section{Conclusion}
Across six matched conditions, fixed partitioning completed all core missions. The short-queue policy exhausted its work before either boat could experience the scheduled slowdown. Reserving twelve tasks at every exchange completed more missions than reserving twelve only after a trend warning. The two exchanges before the first loss of clear-state radio range accounted for all stalls under the always-long policy in this geometry. The two predictive variants showed exceptions.

A common rendezvous rule allowed every policy to complete on fresh seeds. Trend-triggered reservation was faster than reactive short queues in every extension condition. Always-long queues were faster still. Better access to future fade information under the same reservation action did not yield shorter mean missions than always assigning twelve tasks. A prospective evaluation across new geometries and measured channels would establish how far these findings transfer to maritime surveys.
